\begin{table}[t]
\centering
\caption{Accuracy and camera-pair bootstrap intervals by azimuth separation. Positive gain favors CanonFuse3D.}
\label{tab:view_angle}
\small
\setlength{\tabcolsep}{4pt}
\resizebox{\linewidth}{!}{%
\begin{tabular}{lrrrrl}
\toprule
Angle (deg) & Pairs & Canon. avg. & Fused & Mean gain & 95\% CI \\
\midrule
0--30 & 2,160 & 0.2891 & 0.1563 & 0.1328 & [0.1313, 0.1343] \\
30--60 & 2,700 & 0.2888 & 0.1554 & 0.1334 & [0.1321, 0.1347] \\
60--90 & 2,700 & 0.2876 & 0.1544 & 0.1332 & [0.1318, 0.1345] \\
90--120 & 2,700 & 0.2865 & 0.1529 & 0.1337 & [0.1324, 0.1349] \\
120--150 & 2,700 & 0.2852 & 0.1528 & 0.1325 & [0.1314, 0.1336] \\
150--180 & 3,150 & 0.2844 & 0.1532 & 0.1312 & [0.1303, 0.1322] \\
\bottomrule
\end{tabular}}
\vspace{0.3em}
\begin{minipage}{0.97\linewidth}
\footnotesize\textit{Note.} Each camera pair contributes four action-specific sequences; the bins cover all 64,440 test sequences. MPJPE and absolute gain use dataset coordinate units. Intervals resample camera pairs. Detailed paired tests and effect sizes appear in Supplementary Table~\ref{S-tab:view_angle_significance}.
\end{minipage}
\end{table}
